\keywords{Covid-19}
{Agent-based modeling}
{Parallel programming on GPU}
{CUDA}


\pretextualchapter{Abstract}

% O resumo em inglês deve ser organizado em apenas um parágrafo mesmo.
\noindent
Agent-based epidemiological simulations make it possible to represent the evolution of a disease in a population while accounting for the individual and spatial characteristics of the individuals; however, for areas with a large number of individuals and several parameters, they demand high computational power. This work presents the parallelization, using CUDA, of an extended SEIR agent-based model for Covid-19, applied to four Brazilian urban areas: São Paulo, Rocinha, Brasília and Manaus. The GPU implementation is validated against the reference serial implementation under the same parameters, and its performance is assessed through speedup and efficiency. The results show that the parallel implementation reproduces the epidemic curves of the serial implementation, with a root mean square error below 1.2 percentage points, a value compatible with the statistical fluctuation of the Monte Carlo method, and that the speedup grows with the lattice size, reaching about 56 times for the largest city. The performance analysis indicates that the simulation is limited by memory bandwidth, and a data-structure optimization that keeps memory accesses in the cache significantly reduces the execution time.

\printkeys % linha em branco antes
